% 07_vadose.tex - unsaturated columns: RICHARDS, KINEMATIC, LAG and depth classes
\section{Add unsaturated columns}\label{sec:vz}
Under the slopes the water table lies far below the soil: \V{t3:physical:wt7}~m deep on average under HRU~7. With
\cmd{:SoilZoneVadose NONE} (the default), the recharge that leaves the soil reaches the water table on the same day. An
unsaturated column in each cell carries it down instead.

\step{Add the retention curve to the aquifer classes}
The columns need the water retention curve of the aquifer materials: \cmd{POROSITY} (already there), \cmd{CLAPP\_B} and
\cmd{AIR\_ENTRY\_PRESSURE} [m]. Add the two columns to \cmd{:AquiferClasses} in \cmd{Creek.rvp}:
\diff{t3_aquifer.diff}
\why the till values give a drainable porosity close to the till's specific yield. Otherwise Raven warns that a rising
water table would release a different amount of water from the column than MODFLOW stores.

\step{Add RICHARDS columns and run}\label{st:richards}
Add one line at the end of \cmd{Creek.rvg}:
\diff{t3_richards.diff}
\term{t3_run_richards.txt}
\term{t3_summary_rich.txt}
\see \cmd{Exiting Gracefully: Successful Simulation}. The run takes about \V{t3:secrichards}~seconds. There is no new
warning: \cmd{out\_richards/Raven\_errors.txt} holds the \V{t3:nwarnRich} warnings of Step~\ref{st:physrun}. The files of
this step are in \cmd{\$WS/model/step\_10}.

\step{Try KINEMATIC}
Change one word in \cmd{Creek.rvg} and run into \cmd{out\_kinematic}:
\diff{t3_kinematic.diff}
\begin{shell}
$ $RAVEN Creek -o out_kinematic/
\end{shell}
\see \cmd{Exiting Gracefully: Successful Simulation}.

\step{Try LAG}
Change the same word again and run into \cmd{out\_lag}:
\diff{t3_lag.diff}
\begin{shell}
$ $RAVEN Creek -o out_lag/
\end{shell}
\see \cmd{Exiting Gracefully: Successful Simulation}.

\step{Add depth classes}
In a cell on a slope, the DEM points lie at different heights above the water table. With
\cmd{:SoilZoneVadoseClasses 3}, each link between an HRU and a cell gets up to three columns, one per depth class of its
points. Put \cmd{RICHARDS} back and add the line:
\diff{t3_classes.diff}
\begin{shell}
$ $RAVEN Creek -o out_classes/
\end{shell}
\see \cmd{Exiting Gracefully: Successful Simulation}, after about \V{t3:secclasses}~seconds. \cmd{Creek.rvg} is now as in
\cmd{\$WS/model/finished}.
